\chapter{Observable SINR Proxy Derivation}
\label{app:sinr_proxy}

\section{The Observability Constraint}

The true SINR at frame $t$ is defined in Eq.~(\ref{eq:sinr}) as:
\[
    \text{SINR}_t = \frac{(1-\beta_t)P|H_t|^2}{\sigma^2 + \sigma^2_{\text{ICI},t}}
\]
This expression depends on the true channel gain $|H_t|^2$ and the true residual ICI power $\sigma^2_{\text{ICI},t} = C\tilde{\varepsilon}_t^2(1-\beta_t)P|H_t|^2$, neither of which is directly observable at the receiver. Only the received signal $Y_t[k,n]$ is available. The RL state must therefore use proxy quantities computable from $Y_t$.

\section{Step 1: Channel Power Estimation}

After computing the CFO estimate $\hat{\varepsilon}_t$, the receiver applies phase compensation to each received pilot symbol and averages to obtain an estimate of the channel power:
\[
    |\hat{H}_t|^2 = \frac{1}{|\cP_t| N} \sum_{k \in \cP_t} \sum_{n=0}^{N-1} \left|\frac{Y_t[k,n]\, e^{-j2\pi\hat{\varepsilon}_t n T_s}}{\sqrt{\beta_t P / (\rho_t K)}}\right|^2
\]
Dividing by the pilot amplitude $\sqrt{\beta_t P / (\rho_t K)}$ removes the known scaling, leaving a noisy estimate of $|H_t|^2$. The expectation of this quantity under moderate SNR is approximately:
\[
    \mathbb{E}\!\left[|\hat{H}_t|^2\right] \approx |H_t|^2 + \frac{\sigma^2}{\beta_t P / (\rho_t K)}
\]
i.e., the true channel power plus a bias term due to noise. This estimate is used as state variable $\phi_5$.

\section{Step 2: Residual CFO Power Estimation}

The synchronization mismatch metric $\cE_t$ (Eq.~\ref{eq:sync_mismatch}) averages the squared residual between consecutive pilot pairs after CFO compensation. Under the small-angle approximation, when the residual CFO $\tilde{\varepsilon}_t = \varepsilon_t - \hat{\varepsilon}_t$ is small, the dominant contribution to $\cE_t$ is:
\[
    \cE_t \approx (2\pi T_s)^2 \tilde{\varepsilon}_t^2 \cdot \frac{\beta_t P}{\rho_t K} \cdot |H_t|^2 + 2\sigma^2
\]
The term $2\sigma^2$ is the noise floor (two noise samples per pair). Rearranging, an estimate of the residual CFO power $\tilde{\varepsilon}_t^2$ can be extracted:
\[
    \widehat{\tilde{\varepsilon}_t^2} = \frac{\max(0,\;\cE_t - 2\sigma^2)}{(2\pi T_s)^2 \cdot (\beta_t P / |\cP_t|) \cdot |\hat{H}_t|^2}
\]
where $|\cP_t| = \rho_t K$ is the number of pilot subcarriers. The $\max(0, \cdot)$ clamp prevents negative estimates in the noise-dominated regime.

\section{Step 3: Proxy ICI Power}

Substituting $\widehat{\tilde{\varepsilon}_t^2}$ and $|\hat{H}_t|^2$ into the ICI approximation (Eq.~\ref{eq:ici_power}) yields an observable proxy for the ICI power:
\[
    \hat{\sigma}^2_{\text{ICI},t} = C \cdot \widehat{\tilde{\varepsilon}_t^2} \cdot (1-\beta_t)P \cdot |\hat{H}_t|^2
\]

\section{Step 4: Proxy SINR}

The observable proxy SINR used as state variable $\phi_3$ is:
\begin{equation}
    \overline{\text{SINR}}_t = \frac{(1-\beta_t)P\,|\hat{H}_t|^2}{\sigma^2 + \hat{\sigma}^2_{\text{ICI},t}}
    \label{eq:sinr_proxy}
\end{equation}
This quantity is computable entirely from receiver-side observations: $\cE_t$, $|\hat{H}_t|^2$, $\hat{\varepsilon}_t$, and the known system parameters $(\sigma^2, P, \beta_t, \rho_t, T_s, C)$.

\section{Interpretation}

The proxy SINR is an approximation, not the true SINR. Under good synchronization (small $\tilde{\varepsilon}_t$), $\cE_t \approx 2\sigma^2$ and $\widehat{\tilde{\varepsilon}_t^2} \approx 0$, so $\overline{\text{SINR}}_t \approx (1-\beta_t)P|\hat{H}_t|^2/\sigma^2$, the noise-limited SINR. Under poor synchronization, the ICI term inflates the denominator, reducing the proxy SINR. This means $\overline{\text{SINR}}_t$ correctly reflects both the channel quality and the current synchronization state, making it a meaningful state feature for the DQN.

Similarly, the proxy pilot SNR $\overline{\text{SNR}}_{\text{pilot},t} = \beta_t P |\hat{H}_t|^2 / (\rho_t K \sigma^2)$ follows directly from the same channel power estimate and provides complementary information about sensing performance.
